% =====================================================================
% GLOSSARY SOURCE FILE
% Defines every acronym (\newacronym) and term (\newglossaryentry) used in the
% thesis. Acronyms are printed in the Abbreviations list (with page numbers);
% terms are printed in the Terms list.
%
% How to use these in the text:
%   \acrshort{ais}   -> AIS                                   (short form, records the page)
%   \acrlong{ais}    -> Automatic Identification System       (long form)
%   \acrfull{ais}    -> Automatic Identification System (AIS) (long + short)
%   \gls{ais}        -> long form on first use, short form afterwards
%   \acrshortpl{api} -> APIs                                  (plural short form)
%   \gls{anomaly}    -> Anomaly        \glspl{anomaly} -> Anomalies (term, plural)
%
% NOTE: this thesis builds the lists with \makenoidxglossaries + \printnoidxglossary,
% so they are generated by the normal xelatex passes (no makeglossaries step needed).
% =====================================================================

% --- ABBREVIATIONS (acronyms) ---
\newacronym{ais}{AIS}{Automatic Identification System}
\newacronym{api}{API}{Application Programming Interface}
\newacronym{cpu}{CPU}{Central Processing Unit}
\newacronym{crud}{CRUD}{Create, Read, Update, Delete}
\newacronym{geojson}{GeoJSON}{Geographic JavaScript Object Notation}
\newacronym{gpu}{GPU}{Graphics Processing Unit}
\newacronym{http}{HTTP}{Hypertext Transfer Protocol}
\newacronym{id}{ID}{Identifier}
\newacronym{json}{JSON}{JavaScript Object Notation}
\newacronym{mmsi}{MMSI}{Maritime Mobile Service Identity}
\newacronym{nca}{NCA}{Norwegian Coastal Administration}
\newacronym{ntnu}{NTNU}{Norwegian University of Science and Technology}
\newacronym{png}{PNG}{Portable Network Graphics}
\newacronym{rest}{REST}{Representational State Transfer}
\newacronym{ui}{UI}{User Interface}
\newacronym{ux}{UX}{User Experience}
\newacronym{vts}{VTS}{Vessel Traffic Services}
\newacronym{webgl}{WebGL}{Web Graphics Library}

% --- TERMS: core domain objects ---
\newglossaryentry{anomaly}{
    name={Anomaly},
    description={An AIS message flagged as anomalous, for example for an impossible speed, a position jump, or an unexpected maneuver}
}

\newglossaryentry{anomalygroup}{
    name={Anomaly group},
    description={A collection of anomalies grouped together based on shared type and geographic proximity}
}

\newglossaryentry{individualanomaly}{
    name={Individual anomaly},
    description={A single anomaly belonging to an anomaly group, associated with one vessel}
}

\newglossaryentry{basestation}{
    name={Base station},
    description={A shore-based station that receives AIS messages from vessels}
}

% --- TERMS: anomaly types ---
\newglossaryentry{speedanomaly}{
    name={Speed anomaly},
    description={An anomaly where a vessel reports a speed that is physically impossible}
}

\newglossaryentry{jumpinganomaly}{
    name={Jumping anomaly},
    description={An anomaly where a vessel's reported position jumps an impossible distance between consecutive messages}
}

\newglossaryentry{unexpectedmaneuver}{
    name={Unexpected maneuver},
    description={An anomaly where a vessel performs a movement that is implausible for its type or situation}
}

% --- TERMS: anomaly / message data fields ---
% NOTE: short, factual definitions of the variables stored per anomaly / group.
% Verify any details (e.g. exact units, the meaning of "source ID") against the backend schema.
\newglossaryentry{signalstrength}{
    name={Signal strength},
    description={A value representing the strength of the signal of a message between a vessel and a receiver, normalized from 0 to 1 in the dataset}
}

\newglossaryentry{heading}{
    name={Heading},
    description={The compass direction in which a vessel is pointing, given in degrees}
}

\newglossaryentry{lastactivity}{
    name={Last activity},
    description={Timestamp of the most recent message associated with an anomaly group}
}

\newglossaryentry{starttime}{
    name={Start time},
    description={Timestamp of the first anomaly recorded in an anomaly group}
}

\newglossaryentry{location}{
    name={Location},
    description={The geographic position (longitude and latitude) of an anomaly or anomaly group}
}

\newglossaryentry{sourceid}{
    name={Source ID},
    description={Identifier of the source associated with a message, such as the base station or satellite that received it}
    % TODO: confirm exactly what "source ID" refers to in the backend data and adjust if needed.
}

% --- TERMS: AIS message error causes ---
\newglossaryentry{spoofing}{
    name={Spoofing},
    description={The deliberate transmission of false AIS data to misrepresent a vessel's identity, position, or course}
}

\newglossaryentry{jamming}{
    name={Jamming},
    description={Radio interference that disrupts the reception of AIS signals within an area}
}

% --- TERMS: visualization and web-mapping concepts ---
\newglossaryentry{clusteringterm}{
    name={Clustering},
    description={The process of grouping nearby data points based on proximity to reduce visual clutter and improve readability}
}

\newglossaryentry{heatmapterm}{
    name={Heat map},
    description={A visualization method that represents data density through color}
}

\newglossaryentry{scatterplot}{
    name={Scatter plot},
    description={A visualization that plots individual data points on a coordinate plane; a scatter map plots points using longitude and latitude as the axes}
}

\newglossaryentry{webmap}{
    name={Web map},
    description={An interactive map displayed in a web browser}
}

\newglossaryentry{basemap}{
    name={Basemap},
    description={The underlying map layer that data visualizations are drawn on top of}
}

\newglossaryentry{vectortiles}{
    name={Vector tiles},
    description={Geographic data stored in vector form and rendered directly in the browser, used by modern web maps for smoother, customizable rendering}
}

% --- TERMS: data structures ---
\newglossaryentry{geojsonfeature}{
    name={Feature},
    description={A GeoJSON object representing a single geographic entity, consisting of a geometry and a set of properties}
}

\newglossaryentry{featurecollection}{
    name={FeatureCollection},
    description={A GeoJSON object that groups multiple Features into a single collection}
}

% --- TERMS: frameworks and libraries ---
\newglossaryentry{maplibre}{
    name={MapLibre},
    description={An open-source JavaScript library for displaying interactive, vector-tile web maps in the browser}
}

\newglossaryentry{react}{
    name={React},
    description={An open-source JavaScript library for building user interfaces from reusable components}
}

\newglossaryentry{typescript}{
    name={TypeScript},
    description={An open-source, strongly typed superset of JavaScript that adds static type checking and compiles to plain JavaScript}
}

\newglossaryentry{deckgl}{
    name={Deck.gl},
    description={An open-source, WebGL-accelerated JavaScript framework for high-performance visualization of large geospatial datasets}
}

\newglossaryentry{docker}{
    name={Docker},
    description={An open-source platform for packaging and running applications in isolated, portable containers}
}
